% s-239-p239-m002; original PDF 239, 240
\Aphorism{s-239-p239-m002}{332}
Jedes Wort läuft auf Schienen, jeder Begriff wirkt nach ein oder zwei oder wenigen festen Richtungen; mit Wort und Begriff kann man nie eine Curve beschreiben, sondern nur ein Polygon. Durch Verminderung der Sehnenlänge können wir den Anschluss der gebrochenen an die ideale krumme Linie verfeinern; die moderne Sprache, der man erhöhte Biegsamkeit nachrühmt, erreicht diese, indem sie das Wort, den eindeutigen Wortsinn immer kürzer, enger, ärmer nimmt, sodass wir, gegen frühere Zeiten gerechnet, zur Beschreibung desselben Thatbestandes viel mehr Worte brauchen, dabei aber feiner und genauer beschreiben. Die Sprache älterer Zeiten muthet uns an wie ein Netz aus langen, kräftigen Hauptstrichen, die unsere wie ein Geflecht aus lauter kurzen haarfeinen Spinnfäden; jene ist übersichtlicher und energischer, lässt aber vieles Einzelne unberührt; diese überdeckt auch die Zwischenflächen, wirkt aber einigermassen verwirrend.
\par

